\chapter{Introducción}

En este trabajo práctico se aborda el diseño, la implementación en lenguaje SystemVerilog y la validación
funcional de una Unidad de Procesamiento con arquitectura tipo Harvard.

El presente trabajo concentra el desarrollo del punto 4 de la consigna: la integración jerárquica y el análisis
exhaustivo de la unidad procesadora Harvard (\texttt{up\_harvard}), empleando los módulos precedentes como componentes
fundamentales del sistema.

\section{Objetivos del trabajo}

\begin{itemize}
  \item Diseñar en SystemVerilog una unidad de datos modular.
  \item Integrar un banco de registros de tres puertos con una ALU y un Shifter. 
  \item Sintetizar el diseño a fin de inspeccionar el esquemático RTL resultante y corroborar la correspondencia entre
    la descripción en hardware y la topología implementada.
  \item Comprobar el funcionamiento del camino de datos a través de una simulación temporal en la cual se ejecuten
    operaciones de carga, adición, sustracción con constantes y desplazamientos lógicos, registrando las banderas de
    estado generadas.
\end{itemize}

